\documentclass[../main.tex]{subfiles}

\begin{document}
\subsection{\Pointer}

{\Pointer} dapat didefinisikan sebagai suatu variabel yang menyimpan alamat memori. Jika kita mempunyai variabel dengan tipe data tertentu, maka untuk mendapatkan alamat dari variabel tersebut adalah dengan menggunakan operator \citecode{\&} \parencite{raharjo2015pemrograman}.

Kita dapat mendefinisikan sebuah {\pointer} ke tipe data dengan cara mendeklarasikannya dalam bentuk \citecode{*d}, dimana \citecode{d} merupakan nama dari variabel yang didefinisikan. Simbol \citecode{*} harus diulangi untuk setiap variabel {\pointer} \parencite{Lippman_Lajoie_Moo_2013}:

\begin{code}
\begin{verbatim}
int *ip1, *ip2;  // both ip1 and ip2 are pointers to int
double dp, *dp2; // dp2 is a pointer to double; dp is a double
\end{verbatim}
\end{code}

Jika kita mempunyai sebuah variabel yang bertipe \citecode{long} (misalnya \citecode{X}), maka kita dapat memerintahkan {\pointer} \citecode{P} untuk menunjuk ke alamat yang ditempati oleh variabel \citecode{X}. Untuk melakukan hal tersebut kita perlu menuliskan kode seperti berikut \parencite{raharjo2015pemrograman}:

\begin{code}
\begin{verbatim}
// Mendeklarasikan variabel X dengan tipe long
long X;
// Mendeklarasikan pointer P
long *P;

// Memerintahkan P untuk menunjuk alamat
// dari variabel X
P = &X;
\end{verbatim}
\end{code}

\subsection{{\Dereference} dan {\Reference}}
\subsubsection{{\Reference} (\citecode{\&})}
\begin{indentpar}
    Sebuah \textit{reference} merupakan sebuah alias atau sinonim untuk variabel lain. \textit{Reference} dideklarasikan dengan \textit{syntax} \citecode{type\& ref\_name = var\_name;} dimana \citecode{type} merupakan tipe data variabel, \citecode{ref\_name} merupakan nama dari \textit{reference} dan \citecode{var\_name} merupakan nama dari variabel. Sebagai contoh \parencite{hubbard2021programming} :
\end{indentpar}

\begin{code}[n]
\begin{verbatim}
 int& rn = n; // r is a synonym for n
\end{verbatim}
\end{code}
\vspace{1ex}
\begin{indentpar}
    Variabel \citecode{rn} dideklarasikan sebagai sebuah \textit{reference} ke variabel \citecode{n}, yang harus telah terdeklarasi.
\end{indentpar}

\subsubsection{{\Dereference} (\citecode{*})}

\begin{indentpar}
    Nilai sebuah alamat yang ditunjuk oleh variabel {\pointer} diakses dengan menggunakan operator ini \parencite{achmad2018modul}. Operator {\indirect} dengan tanda {\asterisk} (\citecode{*}) ini artinya sebuah variabel yang diawali dengan operator ini diakses secara tidak langsung melalui alamatnya yang disimpan di variabel lain \parencite{munir2007algoritma}.
\end{indentpar}

\begin{code}[n]
\begin{verbatim}
Input = 25;
Fred = Input;
Ted = &Input;
Beth = *Ted;
\end{verbatim}
\end{code}

\end{document}
